\documentclass[11pt,a4paper]{article}
% Scaffold based on the module template; replace bracketed prompts.
% Target 650--800 words including caption and declaration. Compile as you edit.
\usepackage{graphicx}
\usepackage{float}
\usepackage{subcaption}
\usepackage{amsmath}
\usepackage{color}
% Module formatting: do not change.
\usepackage[left=2.5cm,right=2cm,top=1.5cm,bottom=2cm]{geometry}
\renewcommand{\familydefault}{\sfdefault}
\usepackage{minted}
\usepackage{url}
\begin{document}
\begin{center}
{\bf\Large When is a convincing answer enough?}

{\large Benjamin Simpson (SID: 202028988)}\footnote{I used ChatGPT to help interpret the brief, plan the investigation and construct this \LaTeX{} scaffold. [Update this declaration to include any further AI use accurately.]}
\end{center}

\section*{What would convince us?}
% INTRODUCTION: approximately 90--110 words.
[Open with a question about trusting a computed answer. Introduce the
mini-project: implementing Gaussian elimination for $Ax=b$ and comparing it
with NumPy's solver. State what you wanted AI to help with and your initial
expectations. Establish a question about responsible AI use.]

\section*{Putting the answer to the test}
% INVESTIGATION: approximately 200--240 words, plus the caption.
[Describe two or three actual AI interactions: what you requested, what it
provided, and how you checked it. Explain elimination and back substitution
briefly. Introduce partial pivoting as swapping rows to choose the largest
available absolute entry in the current pivot column. Compare your method
with NumPy on the same inputs.]

[Introduce your test family and explain what varies. Choose a reference
solution $x_*$, form $b=Ax_*$ and denote the computed solution by $\hat{x}$.
Distinguish the residual $A\hat{x}-b$ from the solution error $\hat{x}-x_*$.
Mention rounding in forming $b$ if relevant.]

% A short verification snippet; the variables must already be defined.
\begin{minted}{python}
residual = np.linalg.norm(A @ x_hat - b) / np.linalg.norm(b)
error = np.linalg.norm(x_hat - x_star) / np.linalg.norm(x_star)
\end{minted}

[Explain these relative measures, assuming nonzero $b$ and $x_*$, and connect
them to a specific claim you accepted or questioned.]

\begin{figure}[H]
\centering
\IfFileExists{solver_comparison.pdf}{%
\includegraphics[width=0.85\linewidth,height=4cm,keepaspectratio]{solver_comparison.pdf}%
}{%
\fbox{\parbox[c][3.6cm][c]{0.85\linewidth}{\centering
Placeholder: insert your own investigation figure.\\
Suggested: relative solution error against pivot parameter,\\
comparing basic elimination, pivoted elimination and NumPy.}}%
}
\caption{[State the test system, axes, solver labels and measured finding.
Explain any omitted exact-zero errors on logarithmic axes. This placeholder
is not experimental evidence.]}
\label{fig:comparison}
\end{figure}

% Reserve page 2 for reflection. If page 1 overflows, shorten the prose or
% simplify the figure; do not shrink the module's formatting.
\newpage
\section*{An answer is not an explanation}
% DISCUSSION: approximately 240--290 words.
[Interpret the most revealing result in Figure~\ref{fig:comparison}.
Explain how it changed or qualified your judgement of AI's contribution.
Separate algorithmic instability from sensitivity inherent in the problem.
Do not claim that pivoting fixes ill-conditioning or that passing a finite
set of tests proves correctness.]

[Develop one philosophical angle anchored in this experience. Why trust an
established numerical library while questioning generated code? Discuss
assumptions, documentation, testing and accountability. Consider the
counterargument: responsible use does not require reconstructing every
detail of a tool's implementation. Where does understanding become necessary?]

[Compare AI as a code generator, tutor and critic using what you actually
tried. What saved time? What needed extra verification? Could you explain or
modify the method unaided? Identify limitations: one model, a small test
family, or your prior Python experience. Avoid generalising to all students.]

\section*{What I would take into a viva}
% CONCLUSION: approximately 90--120 words, distinct from the discussion.
[State your personal position. Give three concrete recommendations for a
final-year student preparing a report, presentation and viva: for example,
keep an experiment/prompt log, verify important claims independently, and
practise explaining and modifying the method unaided. State what you would
delegate to AI in future and which decisions you would retain. Link these
recommendations to the evidence already discussed.]

% Optional documentation citations as short footnotes near relevant claims:
% https://numpy.org/doc/stable/reference/generated/numpy.linalg.solve.html
% https://netlib.org/lapack/explore-html/d8/da6/group__gesv_ga831ce6a40e7fd16295752d18aed2d541.html
% NumPy calls LAPACK _gesv: LU factorisation with partial pivoting.
% References and the AI declaration must also fit within two pages.
% Final checks: four sections, engaging headings, genuine reflection,
% explained minted snippet, own labelled figure, accurate AI declaration,
% no placeholders or unsupported results, at most two pages.
% This scaffold has not been compiled locally; check pagination in Overleaf.
\end{document}
